\subsection{E(a). Special groups from conjugators and from framed diagrams}
\label{sec:ea}
We give the constructions of \S25, folios 149--154, scans 296--306,
with their precise hypotheses. Let $\pi$ be a discrete hyperbolic
surface group or its full profinite completion, and let $I$ be its finite
peripheral index set. Choose representatives $L_i$ and coherent
parametrizations $c_i:T\to L_i$, where $T=\Z$ or $\widehat\Z$.
The self-centralizing and self-normalizing properties needed below
are the results of Part A. Put
$\mathcal E=\Aut_{\rm lac}(\pi)$ and $\Gamma=\Out_{\rm lac}(\pi)$.
Every $\alpha\in\mathcal E$ has a permutation $\sigma_\alpha$ of $I$
and an automorphism $\chi_\alpha$ of $T$.

\begin{claim}{E.1}{Proved}{The conjugator construction and the product kernel}
Define $S\mathcal E$ to consist of pairs $(\alpha,(g_i)_{i\in I})$
with $\alpha\in\mathcal E$, $g_i\in\pi$, and
\[
 \alpha(c_i(t))=g_i^{-1}c_{\sigma_\alpha(i)}(\chi_\alpha t)g_i
 \quad(t\in T).
\]
The multiplication is
\[
 (\beta,(h_i))(\alpha,(g_i))
 =\bigl(\beta\alpha,
       (h_{\sigma_\alpha(i)}\,\beta(g_i))_{i\in I}\bigr).
\]
There are exact sequences and a canonically embedded normal subgroup
\[
 1\to T^I\to S\mathcal E\to\mathcal E\to1,
 \qquad
 \iota:\pi\hookrightarrow S\mathcal E,
 \quad a\mapsto(\operatorname{Int}(a),(a^{-1})_i).
\]
The groups $T^I$ and $\iota(\pi)$ commute and intersect trivially.
Consequently
\[
 S\Gamma=S\mathcal E/\iota(\pi),\qquad
 1\to T^I\times\pi\to S\mathcal E\to\Gamma\to1,
 \qquad 1\to T^I\to S\Gamma\to\Gamma\to1.
\]
The superscript $!$ takes the inverse image of the kernel of the
permutation action on $I$. All maps in the profinite case are continuous
maps of profinite groups, and all kernels and quotients are closed.
\end{claim}
\begin{proof}
Apply $\beta$ to the defining equation for $\alpha$. The two
conjugating factors multiply as
$\beta(g_i)^{-1}h_{\sigma_\alpha(i)}^{-1}$, which is the inverse of the
factor displayed in the product. Associativity follows either by
substitution in that formula or by associativity of composition of the
specified conjugations. The identity is $(\id,(1)_i)$, and the inverse
has automorphism $\alpha^{-1}$ and $j$th conjugator
$\alpha^{-1}(g_{\sigma_\alpha^{-1}(j)}^{-1})$.

Surjectivity to $\mathcal E$ is the definition of preserving the loop
classes. In its kernel every $g_i$ centralizes $L_i$, hence belongs to
$L_i$; this identifies the kernel with $T^I$. Substituting
$(\operatorname{Int}(a),(a^{-1})_i)$ proves that $\iota$ is a
homomorphism and that conjugation by $(\alpha,(g_i))$ takes $\iota(a)$
to $\iota(\alpha(a))$. Its injectivity follows from $Z(\pi)=1$, including
when $I$ is empty. To check commutation with the kernel, the two products
with $(\id,(\ell_i))$, $\ell_i\in L_i$, both have conjugators
$\ell_i a^{-1}$. Their intersection is trivial because projection to
$\mathcal E$ sends $\iota(a)$ to $\operatorname{Int}(a)$. Finally an
element lying over the identity in $\Gamma$ has an inner automorphism
component and can be multiplied by a unique $\iota(a)^{-1}$ to put that
component equal to the identity. This proves that the whole kernel is
the displayed direct product, not merely an abstract extension of its
two factors.

For a finitely generated profinite $\pi$, its automorphism group has the
usual profinite topology obtained from characteristic open quotients.
The defining equations of $S\mathcal E$ are closed equations in
$\mathcal E\times\pi^I$. Thus it is profinite, the multiplication and
all maps above are continuous, and the compact subgroups giving the
quotients are closed. The same algebraic proof gives the exact
continuous sequences.
\end{proof}

\begin{claim}{E.2}{Proved}{Agreement with autoequivalences of the framed diagram}
The group $S\Gamma$ of E.1 is the group of isomorphism classes of
autoequivalences of
\[
 B_I\longrightarrow\coprod_{i\in I}BL_i\longrightarrow B\pi
\]
with the common orientation action retained. Adding a fixed base object
in $B\pi$ gives $S\mathcal E$. Thus the two constructions on scans
296--300 agree, with the comparison maps and inner ambiguities exactly
as in E.1.
\end{claim}
\begin{proof}
An autoequivalence gives a permutation of the selected objects, the
specified automorphism $\chi$ on their cyclic automorphism groups,
and an automorphism $\alpha$ at the generic base object. The natural
isomorphism making the second square commute is, at its $i$th object,
an element $g_i\in\pi$. Its naturality equation is exactly the defining
equation of E.1. Pasting the squares gives the product there, including
the index $\sigma_\alpha(i)$.

Initially the first square also contains an isomorphism at each selected
peripheral object. A natural transformation at that object changes this
isomorphism by the corresponding element of $L_i$. Use it to make the
first-square isomorphism the identity. After this gauge choice there
is no remaining peripheral-object natural-transformation ambiguity.
A natural transformation at the generic object is an element $a\in\pi$
and changes the representative by $\iota(a)$, with the convention of
E.1. These are precisely the equivalences being divided out in
$S\Gamma$. Fixing the generic base object and its comparison
isomorphism removes that last ambiguity and gives $S\mathcal E$.
Conversely every pair in E.1 defines the three functors and their two
commutativity squares. The two constructions undo each other on the
chosen representatives and on the specified natural transformations,
so the group isomorphisms are canonical relative to the initial
peripheral representatives. Changing those representatives transports
the diagrams by their explicit conjugating isomorphisms.
\end{proof}

\begin{claim}{E.3}{Proved}{The puncture-stabilizer pushout is the normalizer group}
Let $\Gamma_i\subset\Gamma$ stabilize $i$. Restrict the extension
$S\Gamma\to\Gamma$ to $\Gamma_i$ and push its kernel out by
$\operatorname{pr}_i:T^I\to T$. The resulting extension is canonically
\[
 1\to L_i\xrightarrow{\operatorname{Int}}
 \{\alpha\in\Aut_{\rm lac}(\pi):\alpha(L_i)=L_i\}
 \longrightarrow\Gamma_i\to1.
\]
\end{claim}
\begin{proof}
For a representative $(\alpha,(g_j))$ with $\sigma_\alpha(i)=i$, put
\[
 F_i[(\alpha,(g_j))]=\operatorname{Int}(g_i)\alpha.
\]
The defining equation shows that this automorphism normalizes the
actual subgroup $L_i$. Replacing the pair by its generic-basepoint
change replaces $\alpha$ by $\operatorname{Int}(a)\alpha$ and $g_i$ by
$g_i a^{-1}$; the displayed value is unchanged. Hence the map descends
to $S\Gamma_i$. The multiplication formula gives
\[
 \operatorname{Int}(h_i\beta(g_i))\beta\alpha
 =\operatorname{Int}(h_i)\beta\operatorname{Int}(g_i)\alpha,
\]
so it is a homomorphism. An automorphism normalizing $L_i$ has a
representative with $g_i=1$, proving surjectivity. On the kernel $T^I$
the map is $\ell_i\mapsto\operatorname{Int}(\ell_i)$; its kernel is
exactly the other factors because $Z(\pi)=1$. This is the universal
pushout claimed. The normalizer equality in Part A gives the last
kernel $L_i$ in the displayed sequence. This is the compatibility of
scans 305--306, with its actual group of automorphisms rather than an
unspecified outer class.
\end{proof}

\begin{claim}{E.4}{Refuted}{The product-kernel claim requires the rigidity hypotheses}
The formula $\ker(S\mathcal E\to\Gamma)=T^I\times\pi$ does not hold
unchanged for all abelian types. For the unpunctured torus,
$I=\varnothing$, $S\mathcal E=\operatorname{GL}_2(\Z)=\Gamma$ and the
kernel is trivial, whereas $\pi=\Z^2$. For the annulus the inner
basepoint subgroup intersects the peripheral kernel, and the special
outer kernel is a quotient of $T^2$ by one diagonal copy of $T$.
\end{claim}
\begin{proof}
For the torus, all inner automorphisms of $\Z^2$ are the identity and
there are no conjugators indexed by punctures. Thus the pair
construction contains only the automorphism itself, giving the stated
identity map and its trivial kernel.

For the annulus, both peripheral images are the whole group
$\pi=\Z$. The conjugator kernel is $\pi^2$, but the generic-basepoint
change $a$ is $(\id,(a^{-1},a^{-1}))$ and lies in that same kernel.
Dividing by these changes gives $\Z^2/\Delta\Z\cong\Z$, with the
antidiagonal convention after orienting both boundary loops. Keeping
the two end labels and the independent orientation data supplies the
appropriate extended quotient; their omission would create an
additional error. For the plane the peripheral homomorphism
$\Z\to1$ is not injective, so the normalizer calculation used in E.1
is unavailable. Its correct fixed-boundary positive special group is
trivial, as C.20 computes. These explicit maps identify the failed
hypotheses in each case.
\end{proof}
